% =====================================================================
% Appendix: Artifact detail (datasheet, metadata, maintenance)
% -- Relocated from sections/artifacts.tex on 2026-07-26 as part of the
%    structural trim. NOTHING WAS DELETED IN THE MOVE.
% -- Datasheet TODO resolved 2026-07-27 in the \citet{gebru2021datasheets}
%    format. Every required item from the TODO appears below, marked (*).
%
% D&B hygiene. Required in full if NeurIPS Datasets & Benchmarks is the target
% venue; for COLM/ACL the main text keeps one paragraph (sections/artifacts.tex)
% and the detail lives here.
%
% RULE, UNCHANGED: no URL, DOI, or version number may be written here until it
% exists. UPDATED 2026-08-06: the release is now v1.1.0, whose version DOI is
% RESERVED and does not resolve until the deposit is published. v1.0.0 stays
% listed as the superseded record, via \priorversiondoi.
% =====================================================================

\section{Artifact detail}
\label{app:artifacts-detail}

\subsection{Datasheet}
\label{app:artifacts:datasheet}

Following \citet{gebru2021datasheets}.

\paragraph{Motivation.} The artifacts exist to make the paper's own
recommendation checkable against the paper. \S\ref{sec:audit} reports that no
audited source releases per-item outputs; withholding ours would make the
recommendation unfalsifiable. They were created by the author, with no external
funding beyond a university compute allocation.

\paragraph{Composition.} Three kinds of record. (i) \textbf{Per-item outputs}:
88 cell JSONL files, one row per evaluation item, covering every model $\times$
method $\times$ seed $\times$ task in the controlled experiment. (ii)
\textbf{Derived per-cell statistics}: the atlas, at one row per pair-task cell,
with flip and churn quantities, McNemar and TOST outcomes. (iii)
\textbf{Provenance records}: registrations, signed decisions, configs, receipts,
fingerprints and the incident log. \textbf{(*) The atlas contains no
human-subject data and no personally identifying information}; its rows describe
model behaviour on public benchmark items. \textbf{(*) The population is the
public record of compression evaluation, not a census of quantization}
(\S\ref{sec:atlas:caveats}): it is conditioned on what was published and
leaderboard-indexed, so it describes the evidence the field circulates rather
than how quantization behaves in general.

\paragraph{Collection.} S1 cells were mined from public per-item evaluation
dumps; S2 from vendor-published evaluations; the controlled cells were generated
by us inside a pinned container, under protocols frozen before the runs
(\S\ref{sec:prereg}). Pair enumeration was frozen in a manifest before any
statistic was computed.

\paragraph{Preprocessing.} A pair-task cell is admitted only when both sides
carry a binary per-item correctness state and full-prompt hashes match across the
pair. Exclusions are released as a table with a reason per row rather than
summarised, so the admitted population is auditable and not merely asserted.

\paragraph{Distribution and licensing.} Released under \textbf{CC-BY-4.0};
the \texttt{flipeval} package under \textbf{Apache-2.0}.
\textbf{(*) S1's Open LLM Leaderboard v1 archive carries no declared license},
which is recorded as a limitation rather than worked around: we redistribute
\emph{our derived per-cell statistics} and \textbf{not} the raw upstream
per-item files. The frozen pair manifest records the upstream dataset
identifiers and run timestamps per pair, so a third party can re-derive S1 from
the original sources.
% SOURCE for the S1 license limitation:
% docs/ATLAS_MINING_REGISTRATION_2026-07-15.md §2.
\textbf{(*) S2 is Apache-2.0} upstream. Site-specific identifiers (absolute
paths, charge account, usernames, cluster hostnames) were replaced with
placeholders before publication; no accuracy, hash, count, seed, timestamp, job
ID or decision value was altered, and the sealed archives were not rewritten, so
their recorded checksums still verify.
% ADDED 2026-08-02. Final rev-3 checklist §8, Option A. Keep this consistent
% with sections/artifacts.tex; that paragraph is the primary statement and this
% is the datasheet's licensing entry for the same decision.
\textbf{(*) The 17 audited sources are not redistributed.} A redistribution
review on 2026-08-02 found four of them with no third-party republication grant
and the seven method papers under arXiv's default licence, which authorises
arXiv rather than us to distribute them, so the release carries source URLs,
pinned version identifiers, per-file SHA-256 digests, the manifest and a
retrieval script in place of the captured text, and the full-text captures are
kept private (\S\ref{sec:artifacts}). \textbf{(*) Two sources carry provenance
limits no digest can close}, and the manifest records that per row rather than
averaging it away.
% RELOCATED 2026-08-04 (flagship narrative, operation 15) from
% sections/artifacts.tex. The body keeps the fact that two of the seventeen
% bound what a digest can establish and points here; the per-source detail,
% including the expected-drift handling and the fifteen clean re-fetches, is
% carried across whole so nothing was lost in the move.
R11 is served with per-response content: two fetches made seconds apart returned
different bytes, so the digest recorded for it was never a valid fingerprint of
the page, and the retrieval script reports it as expected drift rather than as a
verification failure. R13 does carry a digest for the archived capture, so a
re-fetch can be checked against that capture; what is absent is any digest
recorded \emph{before} the capture was made, so the capture itself was never
corroborated against an independent earlier record. Provenance for these two is
documentary rather than cryptographic, neither is described anywhere in this
paper as hash-verified, and the other fifteen reproduced their recorded digests
exactly when re-fetched on 2026-07-31.

\paragraph{Revision history.} Both revisions are published and the delta is
reported. Rev-1 carried a population defect found by an independent spot-check
that reconciled all 262 compared fields: the arithmetic was right and the
population was not (\S\ref{sec:prereg-spotcheck},
Appendix~\ref{app:prereg:rev2delta}). Rev-2 is what the paper cites; rev-1 is
retained rather than overwritten, because a corrected artifact whose correction
remains visible is this paper's argument applied to itself.

\subsection{Metadata and identifiers}
\label{app:artifacts:metadata}

Released at version \textbf{v1.1.0}, the string shared by the source tag, the
archived release and the dataset revision. The archived DOI is canonical for
citation; the dataset repository is a convenience mirror.

% RESOLVED 2026-07-30 -- see the head of sections/artifacts.tex for how each
% identifier was verified. The version/concept distinction matters here more
% than anywhere: they differ only in the final digit.
% Value column is p{} rather than l: the mirror path used to be hand-split
% across two rows to fit an l column, which a macro cannot do. A p column wraps
% at the \allowbreak points \datasetpath carries.
\begin{tabular}{@{}l p{0.62\linewidth}@{}}
\toprule
Identifier & Value \\
\midrule
% ANONYMITY 2026-07-31: routed through main.tex macros for the TMLR build.
Version DOI (canonical) & \versiondoi \\
Version DOI (v1.0.0, superseded) & \priorversiondoi \\
Concept DOI (latest) & \conceptdoi \\
Source package & \repopath, tag \texttt{v1.1.0} \\
Dataset mirror (secondary) & \datasetpath \\
Container image SHA-256 & \texttt{8260d04c\ldots1db2007} \\
\bottomrule
\end{tabular}

\medskip
\noindent The \textbf{version DOI is what this paper cites}: it resolves to one
frozen state. The concept
DOI is correct only when referring to the artifact series. \textbf{The superseded v1.0.0 DOI is listed
so that a reader holding it can tell which release they have}: its verdict CSV
is rev-2, and \S\ref{sec:artifacts} states what separates the two.
A
\textbf{Croissant record} is served by the dataset repository at its standard
\texttt{/croissant} endpoint and carries the CC-BY-4.0 licence through to the
machine-readable metadata.

\subsection{Maintenance}
\label{app:artifacts:maintenance}

Maintained by the author for \textbf{at least 12 months} from release, via the
source repository's issue tracker, with a target first response of two weeks.
Fixes ship as new tagged versions with a changelog; released versions are never
edited in place, and a superseded revision stays published beside its
replacement. \textbf{Adding atlas pairs requires a dated amendment to the atlas
registration}, not a silent extension. The frozen pair manifest is what makes
the population auditable, so growing it quietly would remove the property the
artifact exists to demonstrate.
